\documentclass[12pt,a4paper]{article}

% Pakete
\usepackage[utf8]{inputenc}
\usepackage[T1]{fontenc}
\usepackage[ngerman]{babel}
\usepackage{mathtools}
\usepackage{amsmath}
\usepackage{amssymb}
\usepackage{amsthm}
\usepackage{geometry}
\usepackage{hyperref}
\usepackage{cleveref}
\usepackage{booktabs}
\usepackage{array}
\usepackage{xcolor}
\usepackage{float}

\geometry{a4paper, margin=2.5cm}

\hypersetup{
  colorlinks=true,
  linkcolor=blue!70!black,
  citecolor=green!60!black,
  urlcolor=blue!60!black,
  pdftitle={Mathematische Analyse der Kegelflaeche},
  pdfauthor={Dogan Balban},
  pdfkeywords={Kegel, Rotationsflaeche, Gauss-Kruemmung, Flaecheninhalt, Volumen}
}

% Theorem-Umgebungen
\theoremstyle{plain}
\newtheorem{satz}{Satz}[section]
\newtheorem{lemma}[satz]{Lemma}
\newtheorem{proposition}[satz]{Proposition}
\newtheorem{korollar}[satz]{Korollar}

\theoremstyle{definition}
\newtheorem{definition}{Definition}[section]
\newtheorem{beispiel}{Beispiel}[section]

\theoremstyle{remark}
\newtheorem{bemerkung}{Bemerkung}[section]

% cleveref
\crefname{satz}{Satz}{Saetze}
\crefname{proposition}{Proposition}{Propositionen}
\crefname{korollar}{Korollar}{Korollare}
\crefname{definition}{Definition}{Definitionen}
\crefname{bemerkung}{Bemerkung}{Bemerkungen}
\crefname{equation}{Gl.}{Gl.}
\crefname{table}{Tab.}{Tab.}
\crefname{section}{Abschnitt}{Abschnitte}

\newcommand{\R}{\mathbb{R}}
\newcommand{\vek}[1]{\vec{#1}}

\title{\textbf{Mathematische Analyse der Kegelfl\"{a}che}\\[0.5em]
\large Vollst\"{a}ndige Kurvendiskussion, Differentialgeometrie\\
und geometrische Eigenschaften}
\author{Dogan Balban}
\date{\today}

\begin{document}
\maketitle

\begin{abstract}
Diese Arbeit pr\"{a}sentiert eine vollst\"{a}ndige mathematische Analyse der Kegelfl\"{a}che,
die durch die Gleichung
\[
z = -\frac{1}{12}\cdot 3^{2} + 6\cdot\sqrt{\frac{x^{2}}{\left(\frac{12}{\pi}\right)^{2}}
+\frac{y^{2}}{\left(\frac{12}{\pi}\right)^{2}}}
\]
gegeben ist. Wir untersuchen algebraische Struktur, geometrische Eigenschaften,
Differentialgeometrie, Kr\"{u}mmungsverhalten und Fl\"{a}chenparametrisierung.
Die Analyse umfasst partielle Ableitungen, Gradient, Normalenvektoren,
Fl\"{a}cheninhalt, Volumen und eine vollst\"{a}ndige Charakterisierung.
\end{abstract}

\tableofcontents
\newpage

\section{Einleitung und Vereinfachung}

\subsection{Gegebener Term}

Wir betrachten die Fl\"{a}che:
\begin{equation}\label{eq:original}
z = -\frac{1}{12}\cdot 3^{2} + 6\cdot\sqrt{\frac{x^{2}}{\left(\frac{12}{\pi}\right)^{2}}
+\frac{y^{2}}{\left(\frac{12}{\pi}\right)^{2}}}
\end{equation}

\subsection{Algebraische Vereinfachung}

\begin{proposition}[Vereinfachung]\label{prop:vereinfachung}
Der gegebene Term \eqref{eq:original} vereinfacht sich zu:
\begin{equation}\label{eq:simplified}
z = -\frac{3}{4} + \frac{\pi}{2}\sqrt{x^2 + y^2}
\end{equation}
\end{proposition}

\begin{proof}
\textbf{Schritt~1} (Konstanter Term):
$-\tfrac{1}{12}\cdot 3^{2} = -\tfrac{9}{12} = -\tfrac{3}{4}$.

\textbf{Schritt~2} (Wurzelterm):
$\dfrac{x^{2}}{(12/\pi)^{2}} = x^{2}\cdot\dfrac{\pi^2}{144}$, analog f\"{u}r $y$.

\textbf{Schritt~3}:
$\sqrt{\dfrac{\pi^2}{144}(x^2+y^2)} = \dfrac{\pi}{12}\sqrt{x^2+y^2}$.

\textbf{Schritt~4}:
$6\cdot\dfrac{\pi}{12}\sqrt{x^2+y^2} = \dfrac{\pi}{2}\sqrt{x^2+y^2}$.

\textbf{Schritt~5}: Gesamtterm: $z = -\tfrac{3}{4}+\tfrac{\pi}{2}\sqrt{x^2+y^2}$.\qed
\end{proof}

\subsection{Darstellung in Zylinderkoordinaten}

\begin{definition}[Zylinderkoordinaten]\label{def:zylinder}
Mit $x = r\cos\theta$, $y = r\sin\theta$, $r = \sqrt{x^2+y^2}$ nimmt
die Gleichung die Form an:
\begin{equation}\label{eq:cylindrical}
\boxed{z = -\frac{3}{4} + \frac{\pi}{2}\,r}
\end{equation}
\end{definition}

\begin{bemerkung}
Gleichung~\eqref{eq:cylindrical} ist unabh\"{a}ngig von $\theta$ -- die Fl\"{a}che ist
eine \textbf{Rotationsfl\"{a}che} um die $z$-Achse, genauer: ein \textbf{gerader Kreiskegel}.
\end{bemerkung}

\section{Geometrische Charakterisierung}

\subsection{Identifikation als Kegel}

\begin{satz}[Kegeleigenschaft]\label{satz:kegel}
Die durch~\eqref{eq:cylindrical} definierte Fl\"{a}che ist ein gerader Kreiskegel mit:
\begin{itemize}
\item Spitze bei $(0,0,-\tfrac{3}{4})$,
\item Achse entlang der positiven $z$-Achse,
\item Steigung $m = \tfrac{\pi}{2}$.
\end{itemize}
\end{satz}

\begin{proof}
Die Gleichung $z = -\tfrac{3}{4}+\tfrac{\pi}{2}r$ ist linear in $r$ mit positivem Koeffizienten.
F\"{u}r $r=0$ gilt $z=-\tfrac{3}{4}$ (Kegelspitze). F\"{u}r $r>0$ w\"{a}chst $z$ linear
-- charakteristisch f\"{u}r einen geraden Kreiskegel.
\end{proof}

\subsection{Kegelwinkel}\label{sec:kegelwinkel}

\begin{definition}[Kegelwinkel]\label{def:kegelwinkel}
Der \emph{Kegelwinkel} $\alpha$ ist der Winkel zwischen der Kegelachse ($z$-Achse)
und einer Mantellinie.
\end{definition}

\begin{proposition}[Kegelwinkel]\label{prop:kegelwinkel}
F\"{u}r unseren Kegel gilt:
\begin{equation}\label{eq:kegelwinkel}
\alpha = \arctan\!\left(\frac{2}{\pi}\right) \approx 0{,}5669\text{ rad} \approx 32{,}48^\circ.
\end{equation}
\end{proposition}

\begin{proof}
Eine Mantellinie verl\"{a}uft von $(r_0,\theta_0,z_0)$ nach $(r_0+\Delta r,\theta_0,z_0+\Delta z)$
mit $\Delta z = \tfrac{\pi}{2}\,\Delta r$. Der Winkel \emph{zwischen Achse und Mantellinie}:
\[
\tan\alpha = \frac{\Delta r}{\Delta z} = \frac{1}{\pi/2} = \frac{2}{\pi}
\quad\Rightarrow\quad
\alpha = \arctan\!\left(\frac{2}{\pi}\right) \approx 32{,}48^\circ.
\]
\end{proof}

\begin{bemerkung}[Elevationswinkel und Komplement\"{a}rwinkel]\label{bem:winkel}
Der komplement\"{a}re \emph{Elevationswinkel} (Winkel zwischen Mantellinie und der
horizontalen $xy$-Ebene) ist:
\[
\beta = \arctan\!\left(\frac{\pi}{2}\right) \approx 1{,}0039\text{ rad} \approx 57{,}52^\circ,
\]
mit $\alpha+\beta = \tfrac{\pi}{2}$ ($32{,}48^\circ+57{,}52^\circ = 90^\circ$).
In einem Teil der Literatur wird $\beta$ als ,,Kegelwinkel'' bezeichnet.
In dieser Arbeit bezeichnet $\alpha$ stets den Winkel zwischen Kegelachse und
Mantellinie (gem\"{a}\ss\ \cref{def:kegelwinkel}).
\end{bemerkung}

\begin{korollar}[\"{O}ffnungswinkel]\label{kor:oeffnung}
Der \"{O}ffnungswinkel (Gesamtwinkel an der Kegelspitze zwischen zwei gegen\"{u}berliegenden
Mantellinien) betr\"{a}gt:
\[
2\alpha = 2\arctan\!\left(\frac{2}{\pi}\right) \approx 64{,}96^\circ.
\]
\end{korollar}

\subsection{Schnittkurven}

\begin{proposition}[Basiskreis]\label{prop:basiskreis}
Der Kegel schneidet $z=0$ im Kreis mit Radius
$r_0 = \tfrac{3}{2\pi} \approx 0{,}4775$.
\end{proposition}

\begin{proposition}[H\"{o}henschnitt]\label{prop:hoehenschnitt}
Der Schnitt mit $z=h$ ($h\geq -\tfrac{3}{4}$) ist ein Kreis mit Radius:
\begin{equation}\label{eq:radials}
r(h) = \frac{2}{\pi}\!\left(h+\frac{3}{4}\right).
\end{equation}
\end{proposition}

\section{Funktionale Darstellung}

\subsection{Explizite Form und erste Ableitungen}

Die Fl\"{a}che als Graph: $f(x,y) = -\tfrac{3}{4}+\tfrac{\pi}{2}\sqrt{x^2+y^2}$,
$(x,y)\in\R^2$.

\begin{proposition}[Partielle Ableitungen]\label{prop:partial}
F\"{u}r $(x,y)\neq(0,0)$:
\begin{align}
f_x &= \frac{\pi x}{2r},\quad
f_y = \frac{\pi y}{2r};\label{eq:grad1}\\
f_{xx} &= \frac{\pi y^2}{2r^3},\quad
f_{yy} = \frac{\pi x^2}{2r^3},\quad
f_{xy} = -\frac{\pi xy}{2r^3}.\label{eq:grad2}
\end{align}
\end{proposition}

\begin{bemerkung}
An der Kegelspitze $(0,0)$ sind die partiellen Ableitungen nicht definiert
(geometrische Singularit\"{a}t).
\end{bemerkung}

\subsection{Gradient und Normalenvektor}

\begin{proposition}[Gradient und Normalenvektor]\label{prop:gradient}
\begin{align*}
\nabla f &= \left(\frac{\pi x}{2r},\;\frac{\pi y}{2r}\right),\quad
|\nabla f| = \frac{\pi}{2}\quad(\text{konstant}),\\
\vek{n} &= \left(\frac{\pi x}{2r},\;\frac{\pi y}{2r},\;-1\right),\quad
\hat{n} = \frac{1}{\sqrt{1+\pi^2/4}}\vek{n}.
\end{align*}
\end{proposition}

\section{Parametrisierung}

\begin{definition}[Parametrisierung]\label{def:param}
\begin{equation}\label{eq:parametrization}
\vek{\gamma}(r,\theta) = \bigl(r\cos\theta,\;r\sin\theta,\;{-\tfrac{3}{4}+\tfrac{\pi}{2}r}\bigr)^\top,
\quad r\geq 0,\;\theta\in[0,2\pi).
\end{equation}
\end{definition}

\begin{proposition}[Tangentialvektoren und Kreuzprodukt]\label{prop:kreuz}
\begin{align*}
\frac{\partial\vek{\gamma}}{\partial r} &= (\cos\theta,\;\sin\theta,\;\tfrac{\pi}{2})^\top,\quad
\frac{\partial\vek{\gamma}}{\partial\theta} = (-r\sin\theta,\;r\cos\theta,\;0)^\top,\\
\frac{\partial\vek{\gamma}}{\partial r}\times\frac{\partial\vek{\gamma}}{\partial\theta}
&= \bigl(-\tfrac{\pi r}{2}\cos\theta,\;-\tfrac{\pi r}{2}\sin\theta,\;r\bigr)^\top,\\
\left|\frac{\partial\vek{\gamma}}{\partial r}\times\frac{\partial\vek{\gamma}}{\partial\theta}\right|
&= r\sqrt{1+\frac{\pi^2}{4}}.
\end{align*}
\end{proposition}

\begin{proof}
Direktes Ausrechnen des Kreuzprodukts; der Betrag:
$\sqrt{(\pi r/2)^2(\cos^2\theta+\sin^2\theta)+r^2} = r\sqrt{1+\pi^2/4}$.
\end{proof}

\section{Fl\"{a}cheninhalt}

\begin{satz}[Mantelfl\"{a}che des Kegelstumpfs]\label{satz:mantelflaeche}
Die Mantelfl\"{a}che des Kegelstumpfs zwischen $r_1$ und $r_2$ ($0\leq r_1<r_2$):
\begin{equation}\label{eq:mantelflaeche}
A = \pi\sqrt{1+\frac{\pi^2}{4}}\cdot(r_2^2 - r_1^2).
\end{equation}
\end{satz}

\begin{proof}
$A = \int_0^{2\pi}\!\int_{r_1}^{r_2}\sqrt{1+\pi^2/4}\cdot r\,dr\,d\theta
= \sqrt{1+\pi^2/4}\cdot 2\pi\cdot\tfrac{r_2^2-r_1^2}{2}
= \pi\sqrt{1+\pi^2/4}\cdot(r_2^2-r_1^2)$.
\end{proof}

\begin{korollar}[Mantelfl\"{a}che von der Spitze]\label{kor:mantelflaeche}
\begin{equation}\label{eq:mantelspitze}
A(r) = \pi r^2\sqrt{1+\frac{\pi^2}{4}} \approx 1{,}8621\cdot\pi r^2.
\end{equation}
\end{korollar}

\section{Volumen}

\begin{satz}[Volumen des Kegelstumpfs]\label{satz:volumen}
\begin{equation}\label{eq:volumen}
V = \frac{4}{3\pi}\left[\left(h_2+\frac{3}{4}\right)^3 - \left(h_1+\frac{3}{4}\right)^3\right],
\quad -\tfrac{3}{4}\leq h_1 < h_2.
\end{equation}
\end{satz}

\begin{proof}
Mit $r(h) = \tfrac{2}{\pi}(h+\tfrac{3}{4})$:
\[
V = \int_{h_1}^{h_2}\pi\,[r(h)]^2\,dh
= \frac{4}{\pi}\int_{h_1}^{h_2}\!\!\left(h+\frac{3}{4}\right)^2 dh
= \frac{4}{3\pi}\!\left[\left(h+\frac{3}{4}\right)^3\right]_{h_1}^{h_2}.
\]
\textit{Kontrolle}: F\"{u}r $h_1=-\tfrac{3}{4}$, $h_2=h$ liefert die Kegelformel
$V = \tfrac{1}{3}\pi R^2 H$ mit $R = r(h)$, $H = h+\tfrac{3}{4}$ dasselbe Ergebnis.
\end{proof}

\section{Kr\"{u}mmungsanalyse}

\subsection{Gau\ss-Kr\"{u}mmung}

\begin{satz}[Gau\ss-Kr\"{u}mmung]\label{satz:gauss}
$K = 0$ f\"{u}r alle regul\"{a}ren Punkte des Kegels.
\end{satz}

\begin{proof}
Kegel sind abwickelbare Fl\"{a}chen (lokal isometrisch zur Ebene): $K=0$.
Direkt: $\kappa_1 = 0$ entlang der Mantellinien (Geraden), also $K=\kappa_1\kappa_2=0$.
\end{proof}

\subsection{Mittlere Kr\"{u}mmung}

\begin{proposition}[Mittlere Kr\"{u}mmung]\label{prop:mittlere}
An einem Punkt mit Radialkoordinate $r>0$:
\begin{equation}\label{eq:mittlere}
H = \frac{\pi}{4r\sqrt{1+\dfrac{\pi^2}{4}}} = \frac{\pi}{2r\sqrt{4+\pi^2}}.
\end{equation}
\end{proposition}

\begin{proof}
F\"{u}r $z = f(x,y)$ gilt:
\[
H = \frac{(1+f_y^2)f_{xx} - 2f_xf_yf_{xy} + (1+f_x^2)f_{yy}}{2(1+f_x^2+f_y^2)^{3/2}}.
\]
Mit $f_x^2+f_y^2 = \pi^2/4$ (konstant) ergibt der Nenner $2(1+\pi^2/4)^{3/2}$.
Der Z\"{a}hler (nach Einsetzen und Vereinfachen mit \cref{prop:partial}):
\[
\frac{\pi y^2}{2r^3}+\frac{\pi x^2}{2r^3}
+\frac{\pi^3(x^4+y^4+2x^2y^2)}{8r^5}
= \frac{\pi}{2r}+\frac{\pi^3}{8r}
= \frac{\pi(4+\pi^2)}{8r}.
\]
Daher:
\[
H = \frac{\pi(4+\pi^2)/8r}{2[(4+\pi^2)/4]^{3/2}}
= \frac{\pi}{2r\sqrt{4+\pi^2}} = \frac{\pi}{4r\sqrt{1+\pi^2/4}}.\qed
\]
\end{proof}

\begin{bemerkung}
\"{A}quivalent mit parametrischer Formel: $H = m/(2r\sqrt{1+m^2})$ mit $m=\pi/2$
liefert $H = (\pi/2)/(2r\sqrt{1+\pi^2/4}) = \pi/(4r\sqrt{1+\pi^2/4})$.
Da $\kappa_1=0$, folgt die zweite Hauptkr\"{u}mmung
$\kappa_2 = 2H = \pi/(2r\sqrt{1+\pi^2/4})$.
\end{bemerkung}

\section{Geod\"{a}tische Linien}

\begin{satz}[Geod\"{a}ten auf dem Kegel]
Die geod\"{a}tischen Linien auf dem Kegel sind:
\begin{enumerate}
\item \textbf{Mantellinien}: Geraden radial von der Spitze nach au\ss en.
\item \textbf{Gewickelte Geod\"{a}ten}: Kurven, die beim Abwickeln des Kegels zu
Geraden in der Ebene werden.
\end{enumerate}
Der abgewickelte Kegel bildet einen Kreissektor mit \"{O}ffnungswinkel
$2\pi\sin\alpha \approx 2\pi\times 0{,}5365 \approx 3{,}370\text{ rad} \approx 193{,}1^\circ$.
\end{satz}

\section{Numerische Beispiele}

\begin{table}[H]
\centering
\caption{Kegelparameter}\label{tab:kegelparameter}
\begin{tabular}{lcc}
\toprule
Parameter & Exakter Wert & Numerisch \\
\midrule
Spitze ($z$-Koord.) & $-\tfrac{3}{4}$ & $-0{,}750$ \\
Steigung $m$ & $\tfrac{\pi}{2}$ & $1{,}5708$ \\
Kegelwinkel $\alpha$ & $\arctan(2/\pi)$ & $0{,}5669$ rad $= 32{,}48^\circ$ \\
Elevationswinkel $\beta$ & $\arctan(\pi/2)$ & $1{,}0039$ rad $= 57{,}52^\circ$ \\
\"{O}ffnungswinkel $2\alpha$ & $2\arctan(2/\pi)$ & $64{,}96^\circ$ \\
Basisradius ($z=0$) & $\tfrac{3}{2\pi}$ & $0{,}4775$ \\
$\sqrt{1+\pi^2/4}$ & $\tfrac{\sqrt{4+\pi^2}}{2}$ & $1{,}8621$ \\
\bottomrule
\end{tabular}
\end{table}

\begin{table}[H]
\centering
\caption{Mantelfl\"{a}che: $A = \pi\sqrt{1+\pi^2/4}\cdot(r_2^2-r_1^2)$}
\label{tab:mantelflaeche}
\begin{tabular}{ccc}
\toprule
$r_1$ & $r_2$ & $A$ \\
\midrule
$0$ & $5$  & $146{,}249$ \\
$0$ & $10$ & $584{,}995$ \\
$5$ & $10$ & $438{,}746$ \\
$0$ & $20$ & $2339{,}979$ \\
\bottomrule
\end{tabular}
\end{table}

\begin{table}[H]
\centering
\caption{Volumen: $V = \tfrac{4}{3\pi}\bigl[(h_2+\tfrac{3}{4})^3-(h_1+\tfrac{3}{4})^3\bigr]$}
\label{tab:volumen}
\begin{tabular}{ccc}
\toprule
$h_1$ & $h_2$ & $V$ \\
\midrule
$-0{,}75$ & $0$  & $0{,}179$ \\
$0$       & $10$ & $527{,}068$ \\
$-0{,}75$ & $10$ & $527{,}247$ \\
$0$       & $20$ & $3791{,}601$ \\
\bottomrule
\end{tabular}
\end{table}

\section{Zusammenfassung}

\begin{table}[H]
\centering
\caption{\"{U}bersicht aller Eigenschaften}
\label{tab:zusammenfassung}
\begin{tabular}{ll}
\toprule
Eigenschaft & Wert / Formel \\
\midrule
Vereinfachte Form & $z = -\tfrac{3}{4}+\tfrac{\pi}{2}r$ \\
Fl\"{a}chentyp & Gerader Kreiskegel \\
Kegelspitze & $(0,0,-\tfrac{3}{4})$ \\
Steigung & $m = \tfrac{\pi}{2}$ \\
Kegelwinkel $\alpha$ & $\arctan(2/\pi)\approx 32{,}48^\circ$ \\
Elevationswinkel $\beta$ & $\arctan(\pi/2)\approx 57{,}52^\circ$ \\
\"{O}ffnungswinkel $2\alpha$ & $2\arctan(2/\pi)\approx 64{,}96^\circ$ \\
Gau\ss-Kr\"{u}mmung & $K=0$ (abwickelbar) \\
Mittlere Kr\"{u}mmung & $H = \dfrac{\pi}{4r\sqrt{1+\pi^2/4}}$ \\[4pt]
Gradient & $|\nabla f| = \tfrac{\pi}{2}$ \\
Mantelfl\"{a}che & $A = \pi\sqrt{1+\pi^2/4}\cdot(r_2^2-r_1^2)$ \\
Volumen & $V = \tfrac{4}{3\pi}\bigl[(h_2+\tfrac{3}{4})^3-(h_1+\tfrac{3}{4})^3\bigr]$ \\
\bottomrule
\end{tabular}
\end{table}

\section{Fazit}

Diese Arbeit hat eine vollst\"{a}ndige mathematische Analyse der Kegelfl\"{a}che
$z = -\tfrac{3}{4}+\tfrac{\pi}{2}r$ pr\"{a}sentiert:

\begin{enumerate}
\item Der komplexe Ausgangsterm vereinfacht sich zu $z = -\tfrac{3}{4}+\tfrac{\pi}{2}r$.
\item Die Fl\"{a}che ist ein gerader Kreiskegel mit Spitze $(0,0,-\tfrac{3}{4})$, Steigung $\pi/2$.
\item Gau\ss-Kr\"{u}mmung $K=0$ (abwickelbare Fl\"{a}che).
\item Kegelwinkel (Achse--Mantellinie): $\alpha = \arctan(2/\pi)\approx 32{,}48^\circ$;
Elevationswinkel: $\beta = \arctan(\pi/2)\approx 57{,}52^\circ$; \"{O}ffnungswinkel: $2\alpha\approx 64{,}96^\circ$.
\item Gradient konstant: $|\nabla f| = \pi/2$.
\item Mittlere Kr\"{u}mmung: $H = \pi/(4r\sqrt{1+\pi^2/4})$.
\item Geschlossene Formeln f\"{u}r Mantelfl\"{a}che und Volumen.
\end{enumerate}

\appendix


\begin{thebibliography}{9}

\bibitem{docarmo}
M.~P.~do~Carmo,
\textit{Differential Geometry of Curves and Surfaces}, revised ed.
Dover Publications, Mineola, NY, 2016.

\bibitem{struik}
D.~J.~Struik,
\textit{Lectures on Classical Differential Geometry}, 2nd~ed.
Dover Publications, New York, 1988.

\bibitem{gray}
A.~Gray, E.~Abbena und S.~Salamon,
\textit{Modern Differential Geometry of Curves and Surfaces with Mathematica}, 3rd~ed.
Chapman \& Hall/CRC, Boca Raton, FL, 2006.

\bibitem{bronstein}
I.~N.~Bronstein und K.~A.~Semendjaev,
\textit{Taschenbuch der Mathematik}, 10.~Aufl.
Verlag Harri Deutsch, Frankfurt am Main, 2020.

\bibitem{pressley}
A.~Pressley,
\textit{Elementary Differential Geometry}, 2nd~ed.
Springer, London, 2010.

\bibitem{balban_helix}
D.~Balban,
\textit{Differentialgeometrische Analyse der Kegelhelix},
Eigenst\"{a}ndige mathematische Publikation, 2026.

\end{thebibliography}

\end{document}
